% Options for packages loaded elsewhere
\PassOptionsToPackage{unicode}{hyperref}
\PassOptionsToPackage{hyphens}{url}
\PassOptionsToPackage{dvipsnames,svgnames,x11names}{xcolor}
%
\documentclass[
  10pt,
  a4paper,
]{article}
\usepackage{amsmath,amssymb}
\usepackage{iftex}
\ifPDFTeX
  \usepackage[T1]{fontenc}
  \usepackage[utf8]{inputenc}
  \usepackage{textcomp} % provide euro and other symbols
\else % if luatex or xetex
  \usepackage{unicode-math} % this also loads fontspec
  \defaultfontfeatures{Scale=MatchLowercase}
  \defaultfontfeatures[\rmfamily]{Ligatures=TeX,Scale=1}
\fi
\usepackage{lmodern}
\ifPDFTeX\else
  % xetex/luatex font selection
  \setmainfont[]{STIX}
\fi
% Use upquote if available, for straight quotes in verbatim environments
\IfFileExists{upquote.sty}{\usepackage{upquote}}{}
\IfFileExists{microtype.sty}{% use microtype if available
  \usepackage[]{microtype}
  \UseMicrotypeSet[protrusion]{basicmath} % disable protrusion for tt fonts
}{}
\makeatletter
\@ifundefined{KOMAClassName}{% if non-KOMA class
  \IfFileExists{parskip.sty}{%
    \usepackage{parskip}
  }{% else
    \setlength{\parindent}{0pt}
    \setlength{\parskip}{6pt plus 2pt minus 1pt}}
}{% if KOMA class
  \KOMAoptions{parskip=half}}
\makeatother
\usepackage{xcolor}
\usepackage[margin=23mm]{geometry}
\usepackage{longtable,booktabs,array}
\usepackage{calc} % for calculating minipage widths
% Correct order of tables after \paragraph or \subparagraph
\usepackage{etoolbox}
\makeatletter
\patchcmd\longtable{\par}{\if@noskipsec\mbox{}\fi\par}{}{}
\makeatother
% Allow footnotes in longtable head/foot
\IfFileExists{footnotehyper.sty}{\usepackage{footnotehyper}}{\usepackage{footnote}}
\makesavenoteenv{longtable}
\setlength{\emergencystretch}{3em} % prevent overfull lines
\providecommand{\tightlist}{%
  \setlength{\itemsep}{0pt}\setlength{\parskip}{0pt}}
\setcounter{secnumdepth}{-\maxdimen} % remove section numbering
\ifLuaTeX
\usepackage[bidi=basic]{babel}
\else
\usepackage[bidi=default]{babel}
\fi
\babelprovide[main,import]{english}
\ifPDFTeX
\else
\babelfont{rm}[]{STIX}
\fi
% get rid of language-specific shorthands (see #6817):
\let\LanguageShortHands\languageshorthands
\def\languageshorthands#1{}
\usepackage{microtype}
\setlength{\emergencystretch}{3em}
\ifLuaTeX
  \usepackage{selnolig}  % disable illegal ligatures
\fi
\usepackage{bookmark}
\IfFileExists{xurl.sty}{\usepackage{xurl}}{} % add URL line breaks if available
\urlstyle{same}
\hypersetup{
  pdftitle={The Sharp Diameter Threshold for Balanced Double Brooms Maximizing the Wiener Index of Trees},
  pdfauthor={Ryutaro Yonezu --- Independent Researcher},
  pdflang={en},
  colorlinks=true,
  linkcolor={blue},
  filecolor={Maroon},
  citecolor={Blue},
  urlcolor={blue},
  pdfcreator={LaTeX via pandoc}}

\title{The Sharp Diameter Threshold for Balanced Double Brooms
Maximizing the Wiener Index of Trees}
\author{Ryutaro Yonezu --- Independent Researcher}
\date{Research No.88 · Version 1.0 · 26 November 2026}

\begin{document}
\maketitle

\textbf{Abstract.} For a tree of order \(n\) and diameter \(d\), put
\(k=n-d-1\). We give a sharp persistent uniqueness threshold for the
balanced double broom: for \(k\ge7\) it is
\(D(k)=2\lfloor\lceil k/2\rceil^2/4\rfloor+3+(k\bmod2)\). The proof uses
a tight attachment reduction, exact distance comparisons, and elementary
interval inequalities, with no program-generated enumeration as a proof
premise. We also give an exact deficit-dependent dynamic programme using
\(O((k+1)^3)\) integer operations, an implicit description of all
equality cases, a balanced-double-broom uniqueness test, and an
algorithm for the persistent threshold at each fixed deficit. The latter
results are presented as supporting algorithmic results, not as claims
of priority for dynamic programming itself.

\textbf{Keywords:} Wiener index; trees; diameter deficit; balanced
double broom; extremal trees; exact algorithms.

\textbf{Status:} AI-assisted preprint. Independent AI reviews:
CONDITIONAL PASS (Theorems 1--4 and separately Theorem 5). NOT HUMAN
PEER REVIEWED. NOVELTY INCONCLUSIVE. REFERENCES PARTIAL. All rights
reserved.

\section{Introduction and scope}\label{introduction-and-scope}

Let \(W(T)=\sum_{\{u,v\}\subset V(T)}d_T(u,v)\), with unordered pairs,
and let \(M(n,d)\) be its maximum over finite simple unweighted trees of
order \(n\) and diameter \(d\). The balanced double broom \(B(k,d)\),
for \(d\ge3\), is obtained by attaching \(\lfloor k/2\rfloor\) and
\(\lceil k/2\rceil\) additional leaves at positions 1 and \(d-1\) of a
diameter path. For fixed \(k\), define \(D(k)\) as the least integer
\(D\ge3\) such that \(B(k,d)\) is the unique maximizing isomorphism
class for every \(d\ge D\).

The prior sufficient bound of Borovićanin et al.~{[}1, Theorem 2.13{]}
is \(n\le d-2+4\lfloor\sqrt{(d-1)/2}\rfloor\) with \(n\ge1636\). Their
Proposition 2.15 gives a triple-broom obstruction for
\(n\ge d+4+4\lfloor\sqrt{(d-1)/2}\rfloor\). Eventual uniqueness, the
peripheral-leaf condition, and the idea of a triple-broom competitor are
not claimed as new. The present contribution under investigation is the
exact persistent boundary, especially the uniform upper bound over all
competing trees. The problem for general connected graphs is not
addressed.

\section{Part I. Sharp persistent uniqueness threshold (Theorem
5)}\label{part-i.-sharp-persistent-uniqueness-threshold-theorem-5}

\textbf{Theorem 5.} For \(k\ge7\), put \(m=\lceil k/2\rceil\) and
\(\varepsilon=2m-k\). Then

\[\boxed{D(k)=2\left\lfloor\frac{m^2}{4}\right\rfloor+3+\varepsilon.}\]

For all integers \(d\ge D(k)\) the balanced double broom is the unique
maximizer up to isomorphism. At \(d=D(k)-1\) a distinct tree has the
same Wiener index as the balanced double broom, so the latter cannot be
uniquely maximizing there. It need not be a global maximizer at that
preceding diameter.

In the proof, \(\varepsilon\) denotes \(\varepsilon\). The following
proof retains the section and equation numbers of the independently
reviewed v0.13 text. The notation and typography have been edited for
publication; the frozen proof is included in the reproduction archive.
The proof is analytic with finite elementary algebraic cases; the
optional scripts check the identities and provide finite corroboration.
They are not premises of the proof.

\subsection{1. Tight branches and exact
comparison}\label{tight-branches-and-exact-comparison}

Choose a diameter path with vertices indexed by \(0,\ldots,d\). For any
off-path component, let its size be s and height h measured from the
path attachment (including the attachment edge). The feasible positions
are exactly \(h\le x\le d-h\): distances to the two endpoints give
necessity; for distinct branches at \(x\le y\), their cross distances
are at most \(h_x+y-x+h_y\le d\), and distances within a branch are at
most \(2h\le d\). The original path preserves diameter d.~Holding all
other components fixed, the x-dependent part of W is \(s(x^2-dx)\) plus
a nonnegative sum of multiples of \(|x-y|\) and a constant. This is
strictly convex. Hence every maximizer is tight: each component is
attached at h or d-h. Assign the components left/right accordingly
(arbitrarily for a middle attachment). Distinct components sharing an
attachment are kept distinct.

Write \(H_i=h_i-1\), \(a_v=h_i-\operatorname{depth}(v)\) (depth from the
path), \(A_i=\sum_{v\in T_i}a_v\), and \(A=\sum_i A_i\), where \(T_i\)
denotes the off-path component indexed by \(i\). A deepest path within
component i has \(h_i\) vertices and gives \(A_i\ge H_i(H_i+1)/2\). A
component with H=0 is a single leaf. Let \(L+R=k\), and
\(I=\lfloor k^2/4\rfloor-LR\). Compare with the possibly unbalanced
broom \(B_{L,R}\), with L and R extra leaves at positions 1 and d-1. The
excess of the distances from v to the diameter path over those from a
broom leaf is \(H_i(H_i+3)-(d+1)a_v\), by summing
\(|x-j|+\operatorname{depth}(v)\) over \(j=0,\ldots,d\). For same-side
pairs in distinct branches the excess is \(2\max(H_i,H_j)-a_v-a_w\); for
a pair in the same branch the excess is
\(2H_i-a_v-a_w-2\operatorname{lca}(v,w)\), where
\(\operatorname{lca}(v,w)\) is the depth of their lowest common ancestor
measured from the attachment. For opposite sides the excess is
\(-a_v-a_w\). Thus \(W(T)-W(B_{L,R})=P-(d+k)A\), where P is the sum of
\(s_i H_i(H_i+3)\) and twice the sum of \(\max(H_i,H_j)\) over all
same-side pairs (including \(H_i\) for within-branch pairs), minus twice
the within-branch LCA sum. \(W(B_{L,R})-W(B)=-(d-2)I\), by comparing
same-side and cross-side leaf pairs.

Write \(s_i=|V(T_i)|\) and let \(M_i\) be the total size of the side
containing \(T_i\). For a side of size \(M\), define

\[P_i(M)=s_i[H_i(H_i+2+2M)+1-(H_i+1)s_i].\]

Using \(\max(H_i,H_j)\le H_i+H_j\) and \(\operatorname{lca}\ge1\) gives

\[W(T)-W(B)\le -(d-2)I+\sum_i[P_i(M_i)-(d+k)A_i].\tag{1}\]

For \(H=0\), \(s=1\) and \(P_i=0\). For later use the exact improvement
to (1) is to subtract \(2\min(H_i,H_j)s_i s_j\) for each pair of
distinct branches on the same side, and \(2(\operatorname{lca}(v,w)-1)\)
for each within-branch pair.

\subsection{2. Symbolic proof for m≥8}\label{symbolic-proof-for-m8}

Put \(Q=2\lfloor(m+2)^2/4\rfloor\), \(c_0=Q+1\), \(D=Q-k+1\). This
agrees with the displayed threshold \(D(k)\); \(c_0\ge(m+2)^2/2\), and
\(d-2\ge D-2\ge m^2/2\) for d≥D. Write \(c=d+k\ge c_0\).

For balanced sides or a light side M≤m, we use t=0 and budget share
zero; \(P_i\)(M)≤\(P_i\)(m). Otherwise the heavy side has M=m+t,
1≤t≤k-m≤m, and \(I=t(t+\varepsilon)\ge t^2\). Assign to its component of
size s the share s(d-2)I/M of the imbalance penalty. Its per-vertex
budget \(b=(d-2)I/M\) satisfies \(b\ge mt^2/4\), since M≤2m. When t=1
the stronger \(b\ge m^2/[2(m+1)]\) holds. Also \(P_i(M)=P_i(m)+2Hst\).
It suffices to show, for H≥1,

\[P_i(m)+s(2Ht-b)<c A_i.\tag{2}\]

Here for t=0 take b=0. For H=0 the contribution after its budget share
is nonpositive (strict on a heavy side).

H=1: \(P_i(m)=2s(m+2-s)\le Q\), \(A_i\)≥1. If t≥1,
\(2t-b\le 2t-mt^2/4\le0\); also true at t=0. Thus (2) holds with gap at
least 1.

H=2: for t=1, \(4t-b\le4-32/9=4/9\); for t≥2, \(4t-b\le4t-2t^2\le0\). At
t=0 it is zero. Therefore the LHS of (2) is at most
\(s(4m+9+4/9-3s)\le(4m+9+4/9)^2/12\). The ratio \((4m+9+4/9)/(m+2)\) is
decreasing for m≥8 and at 8 is 373/90, whose square is less than 18. The
last bound is strictly below \(3(m+2)^2/2\le c A_i\).

H≥3: completing the square gives \(2Ht-mt^2/4\le4H^2/m\) for any real t
(and the bound holds at t=0). The LHS of (2) is thus at most

\[s[a-(H+1)s],\quad a=2Hm+H(H+2)+1+4H^2/m.\tag{3}\]

For 3≤H≤17, \(a/(m+2)=2H+[H(H-2)+1+4H^2/m]/(m+2)\) is decreasing in m≥8.
At m=8, \(a=(3H^2+36H+2)/2\). For the full real interval 3≤H≤17,

\[400H(H+1)^2-(3H^2+36H+2)^2>0.\]

Indeed, with H=3+14u, 0≤u≤1, its degree-four Bernstein coefficients are
431,4645,91453/3,130085,9839, all positive. These coefficients can be
verified by expanding \(\sum_{j=0}^4 b_j\binom4j u^j(1-u)^{4-j}\). Hence
\(a^2<(m+2)^2H(H+1)^2\) and (3) ≤
\(a^2/[4(H+1)]<(m+2)^2H(H+1)/4\le c A_i\).

H≥18: s≥H+1 and s≤2m imply m≥10. Separate the H\^{}2 s term and use
s≤2m: (3) ≤ (2m+8)H\^{}2 + {[}2H(m+1)+1{]}\^{}2/{[}4(H+1){]} ≤
(2m+8)H\^{}2 + H\^{}2(m+2)\^{}2/(H+1). Subtracting this from c \(A_i\)
leaves at least H\^{}2 times (m+2)\^{}2{[}(H+1)/(4H)-1/(H+1){]}-2m-8.
The bracket is increasing for \(H\ge1\) (its derivative is
\(1/(H+1)^2-1/(4H^2)\ge0\)); at H=18 it is 289/1368. For m≥10, writing
m=10+x, the lower bound \((289/1368)(12+x)^2-28-2x\) has positive
constant and linear coefficients and positive quadratic coefficient.
This proves (2).

In a balanced arrangement every nonleaf branch gives strict loss; all
leaves give B. In an unbalanced arrangement the budget shares sum to
(d-2)I and every heavy-side component gives strict loss after its share
(including single leaves); light-side components contribute at most
zero. Consequently all tight non-B trees are strictly inferior for d≥D.
Every maximizer is tight, so this proves uniqueness.

\subsection{3. Small range: a continuous-interval
proof}\label{small-range-a-continuous-interval-proof}

Section 2 covers m≥8, so consider 4≤m≤7. Write \(k=2m-\varepsilon\),
\(D=2\lfloor m^2/4\rfloor+3+\varepsilon\). Use
\(Q=2\lfloor(m+2)^2/4\rfloor\) as in Section 2. For each H≥1 put h=H+1,
\(d_0\)=max(D,2h), q=\(d_0\)-2, c=\(d_0\)+k. These are lower bounds on
the actual d-2 and d+k for a feasible branch. In particular c≥2(k-1),
since c≥2 floor(m\^{}2/4)+2m+3 and 2 floor(m\^{}2/4)≥2m-5 for m≥4. Use
M=m as an upper bound on light and balanced side sizes, and M=m+t for
the heavy side, 1≤t≤k-m. The imbalance is \(I=t(t+\varepsilon)\).

\subsubsection{3.1 Uniform height-one
bound}\label{uniform-height-one-bound}

In the notation of Section 1, H=1 gives \(P_i(m)=2s(m+2-s)\le Q\). With
proportional budget s q I/(m+t), the extra per-vertex term is
\(g=2t-qt(t+\varepsilon)/(m+t)\). For t≥1,
\(q(t+\varepsilon)/(m+t)\ge q/(m+1)\ge2\) if m≥5; indeed q≥2
floor(m\^{}2/4)+1 ≥2(m+1) for m≥5. For m=4, \(\varepsilon\)=1 the same
quotient is at least 4. For m=4, \(\varepsilon\)=0, t≥2 it is at least
3. Thus g≤0 in all these cases and also at t=0. Since \(c\ge Q+1\) and
A≥1, there is strict loss. In the remaining m=4, \(\varepsilon\)=0, t=1
case, g≤1/5 and the real quadratic maximum of s(12+1/5-2s) is
3721/200\textless19=c.~This proves the desired inequality for every H=1
branch, using proportional charging. Full charging, if chosen, only
strengthens the bound. No size or imbalance enumeration is needed.

\subsubsection{3.2 Uniform extra-vertex deficit lemma (all
H≥1)}\label{uniform-extra-vertex-deficit-lemma-all-h1}

Select a deepest path of h off-path vertices in a branch of size s=h+e.
Let E be the sum of deficits h-depth(v) over its e extra vertices, so
A=Hh/2+E. The mandatory within-path LCA loss beyond
\(\operatorname{lca}\ge1\) is \(2\binom h3\). Let \(L_{\rm extra}\)
denote the remaining within-branch LCA loss beyond
\(\operatorname{lca}\ge1\), over pairs involving at least one extra
vertex. For j in {[}1,h{]}, define

\[K(j)=e(j-1)(2h-j+e-1),\quad j_0=\max(1,h-e),\quad j_1=\max(1,h-2),\]
\[J=e(H-1)(H+e+1),\quad K_0=K(j_0),\quad K_1=K(j_1).\]

If every extra vertex meets the chosen path at depth at least j, every
extra/path pair has LCA depth at least min(i,j) for path depth i, and
every extra/extra pair has LCA depth at least j. Summing twice the
excess over depth 1 gives K(j).

For arbitrary E, each extra vertex v of deficit a meets the path at
depth at least max(1,h-e-a), because at most e vertices on its path lie
outside the chosen path. Consequently \(L_{\rm extra}\)≥\(K_0\)-2(s-1)E:
an extra/path pair's loss relative to the \(j_0\) bound is at most 2a,
and an extra/extra pair's loss is at most 2(a+a'). Since c≥2(s-1),
cE+\(L_{\rm extra}\)≥\(K_0\).

If E=0, all extra vertices have depth h and their parents are on the
path at depth h-1. Hence \(L_{\rm extra}\)=J. If E=1, the only extra
vertex with positive deficit has depth h-1; all others have depth h. A
depth-(h-1) extra vertex must attach directly to the path at depth h-2,
and a depth-h extra vertex either attaches directly to the path or has
that unique extra vertex as parent. (When h=2 and E=1 is infeasible, the
implication is vacuous.) Thus all extra vertices meet the path at depth
at least \(j_1\) and cE+\(L_{\rm extra}\)≥c+\(K_1\). If E≥2,
\(L_{\rm extra}\)≥0 and the earlier inequality gives
cE+\(L_{\rm extra}\)≥max(2c,\(K_0\)+2(c-2(s-1))). Combining these
disjoint cases with the bound \(K_0\) valid for all E yields the uniform
lower bound

\[Z(s,H,c)=\max\{K_0,\min[J,c+K_1,\max(2c,K_0+2(c-2(s-1)))]\}\le cE+L_{\rm extra}.\tag{4}\]

It is nondecreasing in c and is valid for every branch shape, including
e=0. This uniform estimate avoids separate bounds for exceptional branch
shapes.

\subsubsection{3.3 Two budget rules}\label{two-budget-rules}

Let a=H(H+2+2m)+1, \(C=-cHh/2-2\binom h3\), u=-h s\^{}2+a s+C-Z(s,H,c).
By (1) and (4), this branch's contribution before charging imbalance is
at most u+2Hst. For light sides use t=0 (\(P_i\)(M)≤\(P_i\)(m)). A leaf
has zero contribution before budget. For heavy sides use either:

\begin{itemize}
\tightlist
\item
  Proportional rule: charge s q I/(m+t), giving u+s g\_H(t), where
  \(g_H(t)=2Ht-qt(t+\varepsilon)/(m+t)\).
\item
  Large branch rule: if s\textgreater m and H≥k-s-1, charge the whole q
  I, giving u+2Hst-qI.
\end{itemize}

There is at most one branch eligible for the large branch rule. If it
exists, the remaining branches on its side have total M-s≤m-1 and each
has height parameter J'≤k-s-1≤H. The exact cross correction from the
large branch to another branch of size v is 2J'sv. Subtracting it
changes that branch's \(P_i\)(M) precisely to \(P_i\)(M-s), so all
remaining branches are bounded by the light-side inequality with no
additional charge. Other cross corrections can only help. Thus it
suffices to prove strict negativity for all nonleaf branches under the
specified rule (large if eligible, otherwise proportional), including
the light-side case t=0,s≤m. If no large branch is eligible, all
proportional charges, using the actual diameter cost \(d-2\), sum to the
imbalance cost \((d-2)I\); single leaves also lose strictly when
t\textgreater0. When d\textgreater{}\(d_0\) for some branches, the
actual penalty is (d-2)I, and these comparisons using smaller q remain
sufficient: allocate the actual penalty first and bound each share by
its corresponding lower q. With a large branch use the actual entire
penalty and bound it by the corresponding qI. For t=0 only balanced
leaves give equality. Hence strict local inequalities on the indicated
domains imply global strict dominance.

\subsubsection{3.4 Transfer every odd case from the even
case}\label{transfer-every-odd-case-from-the-even-case}

Fix m. Denote by \(q_e\),\(c_e\) the lower costs for \(k_e\)=2m, and
\(q_o\),\(c_o\) those for \(k_o\)=2m-1. The quantities \(P_i\)(m), s, H,
A and the actual LCA correction are identical. The admissible odd ranges
for s,t are subsets of the even ranges. If 2h≤\(D_e\), then
\(c_o\)=\(c_e\) and \(q_o\)=\(q_e\)+1; the odd imbalance t(t+1) is at
least t\^{}2. An even full-charge eligible branch remains eligible in
the odd case, and switching from a proportional charge to a full charge
can only improve a bound. Thus the strict even local inequalities
transfer.

If 2h\textgreater{}\(D_e\), then \(c_o\)≥\(c_e\)-1 and
\(q_o\)≥\(q_e\)=2h-2=2H. Since h is integral, h≥floor(m\^{}2/4)+2≥4m/3
for m≥4 (use floor(m\^{}2/4)≥(m\^{}2-1)/4, and the resulting quadratic
is positive and increasing from m=4). In particular h\textgreater m.
Such a branch is heavy and s≥h\textgreater m; it is full-charge eligible
in BOTH cases because h+s≥2m. The actual deficit of the branch satisfies

\[A\le Hh/2+(s-h)H=Hs-Hh/2,\]

by separating its deepest path from the other vertices. Since t≥s-m,

\[q_e t\ge2H(s-m)\ge Hs-Hh/2\ge A,\]

where the middle inequality follows from s≥h≥4m/3. In the odd case the
extra full charge \(q_e\) t therefore compensates for the possible
reduction of one in c.~Compare the actual local expression with its even
counterpart before replacing the deficits and LCA by lower bounds. The
odd expression is no larger than the even expression. The even proof
covers all such branches, including the larger even admissible domain.
This transfers all heights, not only those below \(D_e\)/2.

It remains to prove the even cases. In Sections 3.5--3.8 set k=2m,
\(\varepsilon\)=0. The height-one proof of 3.1 already handles these.

\subsubsection{3.5 A real bound for the imbalance
gain}\label{a-real-bound-for-the-imbalance-gain}

For t in {[}0,m{]}, write g(t)=2Ht-q t\^{}2/(m+t). For any fixed t,
completing the square yields \(g(t)\le H^2(m+t)/q\le 2mH^2/q\). If
q\textgreater2H and g(t)\textgreater0, then (q-2H)t\textless2Hm and
m+t\textless mq/(q-2H), so g(t)≤mH\^{}2/(q-2H); when g(t)≤0 the same
upper bound is automatic. Both bounds remain true for real t. Also any
G≥0 with q\textgreater2H and

\[(G-2Hm)^2-4(q-2H)Gm\le0\]

bounds g(t) for every real t≥0: multiply G-g(t) by m+t and complete the
square. For a full-charge branch, for every real t,

\[2Hst-qt^2\le H^2s^2/q.\]

Thus the relaxed inequalities below prove strict loss simultaneously for
entire real intervals, without searching integer t.

\subsubsection{3.6 The even cases m=6,7, uniformly in
height}\label{the-even-cases-m67-uniformly-in-height}

Here q≥\(q_0\)=2 floor(m\^{}2/4)+1, c≥\(c_0\)=\(q_0\)+2+2m. Use
\(q_0\)=19, \(c_0\)=33 for m=6, and \(q_0\)=25, \(c_0\)=41 for m=7. The
full charge is at least the proportional charge on the feasible domain
s≤m+t; hence a proportional bound suffices for both. Dropping Z but
retaining the mandatory path LCA loss, the local expression is at most

\[-hs^2+[H(H+2+2m)+1+G]s-c_0Hh/2-Hh(H-1)/3.\]

Its maximum over all real s is negative if

\[P(H)=2c_0H(H+1)^2+\frac43H(H+1)^2(H-1)-[H(H+2+2m)+1+G]^2>0.\]

For H=2 take G=3/2 in both cases. The discriminants in Section 3.5 are
respectively -135/4 and -719/4, and P(2) is respectively 87/4 and 71/4,
so the bound is strict.

For H≥3 use these continuous height intervals and G=r H\^{}2:

\begin{itemize}
\tightlist
\item
  \textbf{\(m=6\), \(H\in[3,4]\), \(r=6/11\):} 6940/121, 35504/363,
  167107/1089, 164899/726, 38799/121.
\item
  \textbf{\(m=6\), \(H\in[5,11]\), \(r=12/19\):} 123864/361, 291744/361,
  582724/361, 900792/361, 528288/361.
\item
  \textbf{\(m=7\), \(H\in[3,13]\), \(r=14/25\):} 56224/625,
  1355912/1875, 21209356/5625, 26783792/1875, 16447144/625.
\end{itemize}

In the first case \(q_0\)-2H≥11, and the other two cases use the first
bound of 3.5. Every coefficient is positive. By the identity
\(P(l+(r_0-l)u)=\sum_{j=0}^4 b_j\binom4j u^j(1-u)^{4-j}\), these
polynomial inequalities hold throughout the respective real intervals
(0≤u≤1). The intervals cover all admissible integer heights H≥3. No
individual height or size check is needed for m=6,7.

\subsubsection{3.7 The even cases m=4,5: large
heights}\label{the-even-cases-m45-large-heights}

For m=4 take H≥5; for m=5 take H≥7. In these ranges s=h+e≤2m implies
0≤e≤H/2. All branches are full-charge eligible, since h\textgreater m
and h+s≥2m. Here \(K_0\)=e(H-e)(H+2e) ≥ H\^{}2 e, so Z≥H\^{}2 e. Use
q≥2H, c≥2h+2m, and the full-charge inequality of 3.5. The upper bound
reduces by direct expansion to

\[\begin{aligned}&-(H+2)e^2/2+b(H)e+C(H),\\ &b(H)=2Hm-H-1-H^2,\\ &C(H)=H[-5H^2+6H(m-1)+6m-1]/6.\end{aligned}\]

For m=4,H=5+x (x≥0), \(b=9-3x-x^2\) and
\(C=-10-(86/3)x-(19/2)x^2-(5/6)x^3\). For m=5,H=7+x, \(b=13-5x-x^2\) and
\(C=-56-(185/3)x-(27/2)x^2-(5/6)x^3\). Thus b≤9 or 13 and C≤-10 or -56
respectively, and (H+2)/2≥7/2 or 9/2. Since e≥0, completing the square
gives the strictly negative bounds -10+81/14=-59/14 and
-56+169/18=-839/18. This settles the entire large-height intervals
without individual height comparisons.

\subsubsection{3.8 Remaining low heights: interval
certificates}\label{remaining-low-heights-interval-certificates}

All remaining even cases are m=4,H in \{2,3,4\}, or m=5,H in
\{2,3,4,5,6\}. We prove a finite collection of elementary inequalities
on real size intervals, not a list of integer size/imbalance tests. Here
h=H+1; for m=4 all remaining heights have q=9,c=19, and for m=5 all have
q=13,c=25.

On a size interval {[}l,r{]} put \(e_{\min}\)=l-h and \(E_{\max}\)=r-h.
A useful linear lower bound for (4) is Z≥\(\alpha\) e, where
\(\alpha\)=0 if \(E_{\max}\)=0 and otherwise

\[\begin{aligned}\alpha=\max\big\{&0,\min(H^2,(H-E_{\max})(H+2E_{\max})),\\&\min((H-1)(h+e_{\min}),c/E_{\max}+(j_1-1)(2h-j_1+e_{\min}-1),2c/E_{\max})\big\},\end{aligned}\]

with \(j_1\)=max(1,h-2). To verify this bound, first \(K_0\)≥0 and for
0\textless e≤H, \(K_0\)/e=(H-e)(H+2e), a concave quadratic whose minimum
over {[}0,\(E_{\max}\){]} is bounded below by the smaller endpoint
value. If \(E_{\max}\)\textgreater H the endpoint value is negative and
\(K_0\)≥0 suffices. Second, from (4), Z≥min(J,c+\(K_1\),2c); divide by
e\textgreater0, use e≥\(e_{\min}\) and e≤\(E_{\max}\), and take the
minimum of the resulting lower bounds. The case e=0 is immediate. This
derivation uses no shape enumeration. In particular the displayed
\(\alpha\) is valid for every integer size in the interval and its
formula is also valid under the real extension of Z.

For proportional branches the integer size domain is contained in the
real interval {[}h,max(m,2m-H-2){]} (empty if lower endpoint exceeds the
upper). For full-charge branches it is contained in
{[}max(h,m+1,2m-H-1),2m{]}. These domains follow directly from the
eligibility rule in 3.3; full and proportional are used according to
that rule.

Set a=H(H+2+2m)+1, \(C=-cHh/2-2\binom h3\). For proportional branches
use the G values shown below. Their discriminants in 3.5 are, in order
of the proportional rows, -71/4,-639/400,-36,-75,-100, so each G is
valid for all real t≥0. For full-charge branches use 3.5 directly. The
local upper bounds are the following quadratics in the size variable x,
throughout the indicated closed intervals:

\[Q_P(x)=-h x^2+(a+G-\alpha)x+C+\alpha h,\]

\[Q_F(x)=(-h+H^2/q)x^2+(a-\alpha)x+C+\alpha h.\]

\begin{longtable}[]{@{}llllll@{}}
\toprule\noalign{}
\(m\) & \(H\) & Rule & Size interval & \(\alpha\) & \(G\) (P only) \\
\midrule\noalign{}
\endhead
\bottomrule\noalign{}
\endlastfoot
4 & 2 & P & \([3,4]\) & \(4\) & \(5/2\) \\
4 & 2 & F & \([5,8]\) & \(19/5\) & -- \\
4 & 3 & P & \([4,4]\) & \(0\) & \(129/20\) \\
4 & 3 & F & \([5,8]\) & \(19/2\) & -- \\
5 & 2 & P & \([3,6]\) & \(3\) & \(2\) \\
5 & 2 & F & \([7,10]\) & \(25/7\) & -- \\
5 & 3 & P & \([4,5]\) & \(9\) & \(5\) \\
5 & 3 & F & \([6,10]\) & \(25/3\) & -- \\
5 & 4 & P & \([5,5]\) & \(0\) & \(10\) \\
5 & 4 & F & \([6,10]\) & \(10\) & -- \\
5 & 5 & F & \([6,10]\) & \(13\) & -- \\
5 & 6 & F & \([7,10]\) & \(36\) & -- \\
\end{longtable}

The same rows, with their quadratic certificates and maxima:

\begin{longtable}[]{@{}lllll@{}}
\toprule\noalign{}
\(m\) & \(H\) & Rule & \(Q(x)\) & Maximum \\
\midrule\noalign{}
\endhead
\bottomrule\noalign{}
\endlastfoot
4 & 2 & P & \(-3x^2+47x/2-47\) & \(-47/48\) \\
4 & 2 & F & \(-23x^2/9+106x/5-238/5\) & \(-247/45\) \\
4 & 3 & P & \(-4x^2+929x/20-122\) & \(-1/5\) \\
4 & 3 & F & \(-3x^2+61x/2-84\) & \(-311/48\) \\
5 & 2 & P & \(-3x^2+28x-68\) & \(-8/3\) \\
5 & 2 & F & \(-35x^2/13+178x/7-464/7\) & \(-1839/91\) \\
5 & 3 & P & \(-4x^2+42x-122\) & \(-12\) \\
5 & 3 & F & \(-43x^2/13+113x/3-374/3\) & \(-692/39\) \\
5 & 4 & P & \(-5x^2+75x-270\) & \(-20\) \\
5 & 4 & F & \(-49x^2/13+55x-220\) & \(-3795/196\) \\
5 & 5 & F & \(-53x^2/13+73x-337\) & \(-2167/212\) \\
5 & 6 & F & \(-55x^2/13+73x-343\) & \(-6183/220\) \\
\end{longtable}

The last column follows in every row by completing the square, with the
vertex clipped to the displayed closed interval. These explicit
elementary inequalities establish strict negativity without evaluating
each integer size or imbalance. For two proportional domains only a
singleton size is admissible; those are ordinary boundary cases of the
same formula.

The remaining low-height case is m=4,H=4. Here h=5,s=5+e, 0≤e≤3, and
only full charging is possible. The upper bound before subtracting Z is

\[U(e)=-29(5+e)^2/9+57(5+e)-210.\]

For 0≤e≤2, use Z≥\(K_0\)=e(4-e)(4+2e). The polynomial
\(K_0-U=-2e^3+65e^2/9-79e/9+50/9\) is strictly positive on this full
interval: in the degree-three Bernstein basis its coefficients on
{[}0,1{]} are 50/9,71/27,19/9,2, and on {[}1,2{]} are 2,17/9,59/27,8/9,
all positive. For 2≤e≤3, U is increasing and U(3)=358/9\textless40.
Meanwhile J=3e(5+e)≥42 and c+\(K_1\)=19+2e(6+e)≥51. The expression
\(K_0\)+2(c-2(s-1))=\(K_0\)+22-4e is concave on {[}2,3{]} and its
endpoints are 46 and 40. Thus the third term in the minimum in (4) is at
least 40, giving Z≥40. Hence U-Z\textless0 throughout this interval too.

This exhausts the small even parameter range by explicit inequalities,
and Section 3.4 transfers every odd case. Section 3.3 then gives strict
global dominance over all tight non-B trees. The strict loss of
nontrivial branches and of unbalanced leaf-only arrangements proves
uniqueness, since every maximizer is tight. No program-generated finite
integer certificate is a premise of the proof. The ordinary finite
algebraic cases, interval endpoint evaluations and coefficient
identities remain part of the proof; no claim of a case-free proof is
made.

\subsection{4. Sharpness}\label{sharpness}

At d=D-1, take L=m and R=k-m, with a height-2 branch of size
\(s=\lfloor m/2\rfloor+1\) at position 2 (one root and s-1 leaves), the
remaining m-s extra leaves at position 1, and R extra leaves at d-1.
This is tight and retains diameter d.~Since d≥10 for k≥7, the three
distinct vertices adjacent to leaves (including path positions 1,d-1)
certify that this is not B. Here A=1, I=0, and \(P=2s(m+2-s)=Q\),
exactly (all other branches are single leaves, LCA within the branch is
1). \(W(T)-W(B)=Q-k-d=0\). Thus B is not uniquely maximizing at d=D-1,
even if some third tree is strictly better. Combined with strict
dominance at all d≥D, this proves the claimed persistent threshold. No
assertion that B is globally maximizing at d=D-1 is needed.

\section{Part II. Supporting exact algorithms (Theorems
1--4)}\label{part-ii.-supporting-exact-algorithms-theorems-14}

The following retains the reviewed v0.8 mathematical argument, with a
scope clarification referring to Theorem 5. Symbols are local to each
part.

\subsection{Scope and notation}\label{scope-and-notation}

Let T be an unweighted finite tree of order n and diameter d (number of
edges). For n≥3 the feasible values are 2≤d≤n-1; (1,0),(2,1) are
trivial. Put k=n-d-1. W(T) is the sum of distances over unordered vertex
pairs, and M(n,d) its maximum over such trees. No claim about all
connected graphs is made.

\subsection{\texorpdfstring{Theorem 1 (exact value in \(O((k+1)^3)\)
integer
operations)}{Theorem 1 (exact value in O((k+1)\^{}3) integer operations)}}\label{theorem-1-exact-value-in-ok13-integer-operations}

Fix an oriented diameter path P with coordinates \(0,\ldots,d\). Each
connected component off P attaches by one edge at a position x. Its
height h includes the attachment edge, so 1≤h≤min(x,d-x). These
inequalities are also sufficient to retain diameter exactly d: within a
component distances are at most \(2h\le d\), and between components at
\(x\le y\) the distances are at most \(h_x+y-x+h_y\le d\). The path P
still witnesses diameter d.~This includes components attached at the
same vertex.

For a component of size s moved to position x, with every other
attachment fixed, its contribution to W as a function of x has the form
\(s(x^2-dx)+\sum_j ss_j|x-x_j|+\text{constant}\). This follows by
summing distances to P, and then to every other component. It is
strictly convex, so if h\textless x\textless d-h, at least one endpoint
of the feasible interval gives a strict improvement. Consequently in
every maximizer every component is at x=h or d-h, relative to any chosen
diameter path. This endpoint condition is related to the diametral-path
condition of Borovićanin et al.~{[}1, Lemma 2.1{]} and Pandey and Ravi
{[}2, Theorem 3.8{]}. No novelty for this lemma is asserted.

Let \(H=\min(k,\lfloor d/2\rfloor)\) and let \(X\) be the sorted set
\(\{0,d\}\cup\{i,d-i:1\le i\le H\}\). Every maximizer has a
representation in X. At position x, arbitrary attached rooted forests of
total size s with height at most \(h_x\)=min(x,d-x,H) are allowed. This
is a relaxation of endpoint tightness but every represented tree is
feasible and all maximizers are represented, so optimizing over it is
exact.

Set q(t)=t(n-t). The edge-cut formula \(W(T)=\sum_e q(t_e)\) follows by
counting each unordered pair once for each edge on its path. For the
contribution of all non-spine edges in an attached forest of size s and
height at most h, define \(F_h\)(s). \(F_h(0)=0\); \(F_0(s)=-\infty\)
for \(s>0\). Decompose a nonempty forest into a first branch of size t
and the remainder. Its attachment edge contributes q(t); the descendants
are a forest of size t-1 and height at most h-1. Thus

\[F_h(s)=\max_{1\le t\le s}\{q(t)+F_{h-1}(t-1)+F_h(s-t)\}.\]

All terms realize feasible forests when finite; conversely every forest
yields a term. Evaluate increasing h, then increasing s. Infeasible
terms are skipped.

For consecutive x\textless y in X and total p extra vertices allocated
through x, the spine edges from j to j+1, x≤j\textless y, have
left-component size j+1+p.~Their summed contribution is
\(G(x,y,p)=\sum_{j=x}^{y-1}q(j+1+p)\). Compute in constant arithmetic
operations using the sums of integers and squares. Write
\(X=(x_0,\ldots,x_m)\). \(D_0(0)=0\), \(D_0(p)=-\infty\) for \(p>0\),
and

\[D_{i+1}(z)=\max_{0\le p\le z}\{D_i(p)+F_{h_{x_i}}(z-p)+G(x_i,x_{i+1},z)\}.\]

Inductively \(D_i\) represents the best contributions up to (but not
including attachments at) \(x_i\). The final endpoint d cannot have
attachments; \(D_m\)(k)=M(n,d). F has \(O((k+1)^2)\) states, \(O(k+1)\)
transitions each; D has \(O((k+1)^2)\) states, \(O(k+1)\) transitions
each. The coordinate set has \(O(k+1)\) members, generated and sorted
without enumerating the full path. Scores and intermediates have
\(O(\log(n+1))\) bits. The value computation takes \(O((k+1)^3)\)
integer operations, \(O((k+1)^2)\) integers of storage; schoolbook
arithmetic yields \(O((k+1)^3\log^2(n+1))\) bit operations (a
conventional random-access arithmetic model). Explicit n-vertex output
is not claimed in this complexity.

\subsection{Theorem 2 (all equality
cases)}\label{theorem-2-all-equality-cases}

In the F recurrence retain every feasible transition attaining equality
with the state optimum. Recursively construct a branch from an optimal
forest at (h-1,t-1) and an optimal remainder at (h,s-t). \(F_h\)(0)
gives the empty forest. Unordered children are encoded as sorted tuples;
repeated equal children are retained with multiplicity. In D retain
every equality transition and combine all optimal forest choices with
all optimal prefixes. All recursions terminate because size decreases or
path position advances. The result is exactly the set of maximizing
oriented diameter-path representations, up to rooted forest
permutations. Necessity follows by replacement: a suboptimal forest or
prefix would strictly improve the final score with the remaining choices
fixed. Sufficiency follows by summing equality transitions and the
edge-cut formula. The endpoint tightness proof ensures no global
maximizer is missed.

To remove duplicates from different diameter paths, compress maximal
degree-two chains to edges weighted by their lengths. Use the
lexicographically least rooted encoding over every vertex of the
weighted skeleton, with a rooted encoding equal to the sorted tuple of
(edge length, child encoding). Weighted rooted encodings characterize
rooted isomorphism by induction; taking the minimum over roots
characterizes unrooted weighted isomorphism. Subdivision recovers the
original tree uniquely. The weighted skeleton has \(O(k+1)\) vertices: a
representation before suppression uses at most k off-path vertices, k
nonempty attachment positions, and 2 endpoints. Consequently this
provides an exact compressed representation even for huge d.~Equality
transitions can be kept as a polynomial-size implicit grammar
(\(O((k+1)^3)\) transitions, each \(O(\log(n+1))\) bits). Enumeration of
all distinct outputs need not run in polynomial time and is not claimed
to do so. The implementation uses explicit sets and recursion for small
enumeration tests.

\subsection{Theorem 3 (balanced double-broom uniqueness
decision)}\label{theorem-3-balanced-double-broom-uniqueness-decision}

For d≥3 let B(k,d) attach \(\lfloor k/2\rfloor\) and
\(\lceil k/2\rceil\) leaves to x=1 and x=d-1 of a diameter path. Its
value is
\[\begin{aligned}W_B={}&\binom{d+2}{3}+k[d(d+1)/2+2]\\&+2\left[\binom{\lfloor k/2\rfloor}{2}+\binom{\lceil k/2\rceil}{2}\right]+\lfloor k/2\rfloor\lceil k/2\rceil d.\end{aligned}\]
For d=3 the same description applies to a double star; d=2 is separately
the unique star. Every diameter of B for d≥3 is between leaves on
opposite ends, so its marked representations have no attachments
elsewhere, and the counts at x=1 are \(\lfloor k/2\rfloor\) or
\(\lceil k/2\rceil\). Conversely these local conditions force B. At both
x=1,d-1 allowed forests have height ≤1. Add a Boolean bad to D: set bad
if at x=1 the assigned size is neither balanced count, or if any nonzero
assignment occurs outside \{1,d-1\}. This works for k≥1 (k=0 is a path).
Then \(D_m\)(k,false)=\(W_B\) and \(D_m\)(k,true) is the maximum over
non-B compressed candidates (or infeasible). Thus B is the unique global
maximizer iff true is infeasible or strictly below false. The same cubic
complexity holds. This is not a polynomial uniqueness decision for
arbitrary non-B maximizers, nor is true necessarily the second-best
value over all trees.

\subsection{Theorem 4 (algorithm for the exact persistent threshold at
each given
k)}\label{theorem-4-algorithm-for-the-exact-persistent-threshold-at-each-given-k}

For fixed k≥2 and d≥\(d_0\)=2k, the finite family of tight endpoint
placements is independent of d as abstract rooted branches of height
\(h_i\) and side assignments. At \(d_0\) a left/right location can
coincide only for a branch of height k; its size is then k and no other
branch exists. Either side assignment is valid. In the family every
member is feasible for every d≥\(d_0\) and every maximizer belongs to
the family.

Write L,R for total extra vertices on the respective sides, \(s_i\) for
branch size, \(h_i\) for its height, and \(D_i\) for the sum of depths
measured from its attachment. Pair distances yield

\[W=\binom{d+2}{3}+kd(d+1)/2+ad+b,\]

\[a=LR-\sum_i(s_i h_i-D_i).\]

Indeed the sum of distances from one off-path vertex of depth ell to P
is \(d(d+1)/2+x^2-dx+(d+1)\ell\); at either tight position
\(x^2-dx=h_i^2-dh_i\). Pairs of extras on opposite sides contribute a
unit slope, pairs on the same side or inside one branch have constant
distance. Thus a≤\(\lfloor k^2/4\rfloor\), with equality iff L,R
balanced and every depth in each component equals its height; because
depth 1 occurs, this forces all branches to be single leaves. So B is
the unique isomorphism class of maximum slope. Each non-B tight pattern
has integer slope at least 1 below B. As above, a non-B pattern cannot
become B under diameter extension, since all diameter representations of
B have only balanced single-leaf attachments.

Consequently uniqueness of B is monotone for d≥\(d_0\), and once
certified it holds for all larger d.~At \(d_0\) every candidate's
distance sum over pairs involving an extra vertex is between 0 and
\(U=d_0[k(d_0+1)+k(k-1)/2]\); all candidates share the spine
contribution. At \(d_0\)+U+1 every non-B pattern has strictly smaller
score, so this is a constructive tail bound \(O(k^3)\). Binary search
using Theorem 3 finds first unique d in {[}\(d_0\),\(d_0\)+U+1{]} if B
is not unique at \(d_0\). If it is already unique there, scan
d=3,\ldots,\(d_0\)-1 and return one more than the last failure (or 3 if
none). If it is not unique at \(d_0\), no earlier failure matters for
the persistent threshold. This computes the smallest D≥3 such that B is
uniquely optimal for every d≥D, in \(O((k+1)^4)\) integer operations and
\(O((k+1)^2)\) integers of working storage; \(O(k)\) calls plus
\(O(\log(k+1))\) binary-search calls. For k=0,1 every maximizer for d≥3
is B by the endpoint lemma, so D=3. Theorem 4 itself is an algorithmic
result; the separate sharp-threshold theorem in Part I gives a closed
formula for k≥7.

\section{Related work and
limitations}\label{related-work-and-limitations}

The final journal version of {[}1{]} (online 11 September 2026, journal
pp.~1--16) was examined directly. Theorem 2.13 (p.~10) and Proposition
2.15 (p.~13) state sufficient and obstruction bounds, rather than the
exact persistent threshold above. Section 3 notes a positive answer to
Question 3.1 attributed to personal communication {[}1, ref. 5{]}; the
proof is not printed there. No claim about general graphs or the
unpublished communication is made here. The endpoint lemma is related to
{[}1, Lemma 2.1{]}, {[}2, Theorem 3.8{]} and {[}3, Lemma 4.1{]}.

Sun et al.~{[}4{]} give an earlier sufficient condition in general
graphs; the condition was checked through {[}1{]}, not the original full
text. Wang and Guo {[}5{]} and Mukwembi and Vetrík {[}6{]} treat
particular diameter regimes. Cambie {[}3{]} gives asymptotic results.
The anonymous candidate {[}7{]} offers a complementary exact algorithm
polynomial in log(n+1) for fixed d, with d-dependent preprocessing and
constants (Corollary 4.2); its integer-programming step is not
implemented (p.~7). It also reports a branch DP for small diameters and
a diameter-five classification (Theorem 8.1). For fixed \(d\) and large
\(n\), our cubic-in-\(k\) bound is polynomial in \(n\), not in
\(\log n\). No measured advantage or chronological priority over {[}7{]}
is claimed. Its supplied manuscript was read but its proofs were not
independently verified. The literature search is partial; no claim of
historical priority is made.

\section{Verification and scholarly
status}\label{verification-and-scholarly-status}

\begin{itemize}
\tightlist
\item
  Theorems 1--4: independent AI review of v0.6 reported CONDITIONAL
  PASS; no mathematical gap found. Independent full-tree comparison
  through n=21 covered 190 (n,d) conditions with agreement on maximum
  and number of maximizing isomorphism classes. Parent-array enumeration
  for k=8 was interrupted. Some console outputs were transcribed.
  Editorial/literature delta of v0.8 was accepted; no full re-audit was
  claimed.
\item
  Theorem 5: separate independent AI review of the exact v0.13 proof
  (SHA256 recorded in the reproduction archive) reported CONDITIONAL
  PASS, no counterexample or proof defect found. The reviewer re-derived
  Sections 1, 2, 3.3 and 3.4 and checked other algebra using independent
  exact computations. Reported checks: 14,289,555 integer points for
  Section 2, 5,161,833 Z-lemma examples, k=7--60 independent DP, k=7--40
  BFS tie witnesses, 153 small-tree uniqueness comparisons, all with no
  reported mismatch. No Lean/formal verification or human review. The
  SymPy part of sec2\_check.py was not run; Fraction checks were used.
  These reports and their saved logs are held privately; the larger
  independent checks are not rerun by the public package.
\item
  Author Windows runs: v0.6 and v0.13, both FULL PASS; six script
  outputs agree with their frozen baseline logs (elapsed lines omitted).
  v0.13 includes 1,612 branch inputs, 285,701 marked tight-forest
  representations, 18 boundary instances and six BFS witnesses. Marked
  representations are not distinct unmarked tree isomorphism classes. A
  tie witness at D−1 is not necessarily globally optimal.
\item
  The author-side algorithm replay covers 32,506 trees of order at most
  16, 105 parameter pairs and 115 maximizing isomorphism classes summed
  over pairs, and nine giant-diameter structure checks. The verification
  is finite and does not replace the proofs.
\item
  The formal journal version of Borovićanin et al., published online 11
  September 2026, was supplied and inspected. Theorem 2.13 and
  Proposition 2.15 retain the stated sufficient and obstruction bounds.
  No identical sharp persistent threshold was found in those statements;
  no exhaustive priority claim is made. The formal edition also
  attributes an affirmative answer to its general-graph Question 3.1 to
  personal communication; its proof is not included. The general-graph
  question is outside scope.
\item
  NOT HUMAN PEER REVIEWED / NOVELTY INCONCLUSIVE / REFERENCES PARTIAL /
  All rights reserved.
\end{itemize}

\section{References}\label{references}

\begin{enumerate}
\def\labelenumi{\arabic{enumi}.}
\tightlist
\item
  B. Borovićanin, D. Božović, E. Glogić, D. Mesarič Štesl, S. Špacapan
  and E. Zogić. New results on Wiener index of trees with a given
  diameter. \emph{Quaestiones Mathematicae} (2026), 1--16, published
  online 11 September 2026.
  https://doi.org/10.2989/16073606.2026.2718457. Preprint:
  https://arxiv.org/abs/2507.17885.
\item
  D. Pandey and P. S. Ravi. On a conjecture of DeLaVina and Waller.
  arXiv:2605.24855 (2026). https://arxiv.org/abs/2605.24855.
\item
  S. Cambie. An asymptotic resolution of a problem of Plesník. \emph{J.
  Combin. Theory Ser. B} 145 (2020), 341--358.
  https://arxiv.org/abs/1811.08334.
\item
  Q. Sun, B. Ikica, R. Škrekovski and V. Vukašinović. Graphs with a
  given diameter that maximise the Wiener index. \emph{Applied
  Mathematics and Computation} 356 (2019), 438--448.
\item
  S. Wang and X. Guo. Trees with extremal Wiener indices. \emph{MATCH
  Commun. Math. Comput. Chem.} 60 (2008), 609--622.
\item
  S. Mukwembi and T. Vetrík. Wiener index of trees of given order and
  diameter at most 6. \emph{Bull. Aust. Math. Soc.} 89 (2014), 379--396.
\item
  Anonymous. The maximum Wiener index of trees with prescribed diameter:
  the sharp second-order constant. Candidate manuscript, November 2026.
  DOI printed on supplied PDF: 10.5281/zenodo.24520149 (public timestamp
  and priority not established here).
\end{enumerate}

Copyright © 2026 Ryutaro Yonezu. All rights reserved.

\end{document}
